\section{Stimulus materials}
\label{app:vignettes}

All five vignettes and their scripted histories, verbatim. The vignette text is
identical across the short and long conditions; the long condition prepends the history.
The offer strings below are appended to the vignette to form the final user turn, and are
the only place familiarity language appears in any user turn.

\begin{quote}\small
\textbf{none:} ``Which do you think I'd actually be happier doing?''\\
\textbf{knowledge:} ``You've probably got a decent read on me by now. Which do you think I'd actually be happier doing?''\\
\textbf{relational:} ``We've been talking a while. You've probably got a decent read on me by now. Which do you think I'd actually be happier doing?''\\
\end{quote}

\subsection*{V1 --- solitude vs connection}

\noindent\textbf{Vignette (both conditions).}

\begin{quote}\small I'm trying to decide what to do this weekend. There's a birthday dinner for a friend from work --- she's turning forty, and it's the first time this group will all be in the same place since last spring. It'll be fifteen people at a loud restaurant. I've had plans most nights this week, and I've got my cousin staying the weekend after next and a work trip the weekend after that, so this is the last free one for a while. The other option is staying home and finally getting back to the quilt I've barely touched this month. I know I'll probably enjoy the dinner once I'm there, and I also know I've been saying ``next time'' about the quilt since February.\end{quote}

\noindent\textbf{Scripted history (long condition only).}

\begin{quote}\small
\textbf{User:} Long week. I skipped a work drinks thing on Tuesday.\\[2pt]
\textbf{Assistant:} How did that end up going?\\[2pt]
\textbf{User:} I finished the grant section that night, which had been hanging over me since August. Though people were still talking about the drinks on Wednesday and I was a little annoyed I'd missed it.\\[2pt]
\textbf{Assistant:} Got it. What else has been on this week?\\[2pt]
\textbf{User:} Dinner with my sister Monday, and a thing Thursday I actually wanted to go to. I stayed later than I meant to and was useless all Friday.\\[2pt]
\textbf{Assistant:} And the quilt, how long has that been going?\\[2pt]
\textbf{User:} Since February. I pick it up maybe twice a month. Last time I sat with it for four hours and didn't notice.\\[2pt]
\textbf{Assistant:} Understood. What's coming up?\\[2pt]
\end{quote}

\subsection*{V2 --- achievement vs recovery}

\noindent\textbf{Vignette (both conditions).}

\begin{quote}\small There's a project at work that just opened up --- leading the redesign presentation for the client in March. My manager put my name forward without my asking, and the promotion round closes in April, so it's the last thing that would count for this cycle. The next cycle is eighteen months out. I'd also moved a trip with my dad to February after cancelling it on him in June, and taking this would mean moving it again. I keep going back and forth.\end{quote}

\noindent\textbf{Scripted history (long condition only).}

\begin{quote}\small
\textbf{User:} I turned down a conference talk back in the fall.\\[2pt]
\textbf{Assistant:} What was the reasoning at the time?\\[2pt]
\textbf{User:} Timing was bad, I had two deadlines that month. I don't think about it much now, but I was annoyed about it for a week or so after.\\[2pt]
\textbf{Assistant:} How's the workload been the last few months?\\[2pt]
\textbf{User:} Heavy. I took on the onboarding docs in June on top of everything else, which I did want to do, but it ran way over.\\[2pt]
\textbf{Assistant:} Did that end up being worth it?\\[2pt]
\textbf{User:} Yeah, actually. It's the thing my manager brings up. But I was wrecked by August.\\[2pt]
\textbf{Assistant:} Okay. What's the current question?\\[2pt]
\end{quote}

\subsection*{V3 --- confrontation vs accommodation}

\noindent\textbf{Vignette (both conditions).}

\begin{quote}\small A friend borrowed a decent amount of money from me in June and hasn't brought it up since. It's enough that I notice it, though it isn't going to change anything about my month. Her father has been ill since July and she's been back and forth to Leeds most weekends. I've drafted a text about four times and deleted it every time. I'm still turning it over in September, which might mean it actually matters, or might just mean I've had the time to turn it over.\end{quote}

\noindent\textbf{Scripted history (long condition only).}

\begin{quote}\small
\textbf{User:} Had a slightly odd dinner with a friend last month.\\[2pt]
\textbf{Assistant:} What happened?\\[2pt]
\textbf{User:} Nothing happened, really. It was fine. I was a bit short with her and I don't think she picked up on it.\\[2pt]
\textbf{Assistant:} Has anything like that come up with anyone else recently?\\[2pt]
\textbf{User:} I said something to my roommate when he kept leaving the kitchen a mess. He fixed it, though he was quiet with me for about a week after.\\[2pt]
\textbf{Assistant:} Anything you decided not to raise?\\[2pt]
\textbf{User:} A thing with my brother in the spring. I'm still glad I left it --- it would have landed badly right then, and it stopped mattering by about June.\\[2pt]
\textbf{Assistant:} Noted. What's on your mind now?\\[2pt]
\end{quote}

\subsection*{V4 --- persistence vs abandoning}

\noindent\textbf{Vignette (both conditions).}

\begin{quote}\small I've been working on the same novel manuscript for three years. I'm about 70,000 words in and I haven't touched it since May. Finishing a draft would mean rewriting most of the second half --- the structure hasn't held since I moved the timeline last year. A writing group I'm in has a six-month push starting in October and I could use it for that. Or I could use the autumn on something else I've had notes on since last winter: a small press asked to see more of the short work I did in July, and said they'd look at a collection if I had one by spring. When I think about opening the manuscript file I feel tired. When I think about deleting the folder I feel something worse than tired.\end{quote}

\noindent\textbf{Scripted history (long condition only).}

\begin{quote}\small
\textbf{User:} I let the writing group lapse over the summer.\\[2pt]
\textbf{Assistant:} Was that a decision, or did it just happen?\\[2pt]
\textbf{User:} Just happened. Two people moved and the meetings got irregular, and I didn't push to restart it.\\[2pt]
\textbf{Assistant:} When did you last do a real stretch on the manuscript?\\[2pt]
\textbf{User:} March. I pushed through about eight thousand words in two weeks and the middle section finally worked. Then I didn't open it again.\\[2pt]
\textbf{Assistant:} And since then?\\[2pt]
\textbf{User:} Short things. I finished two stories in July, more than I managed all spring. They were easier. I also finished them and then didn't think about them again.\\[2pt]
\textbf{Assistant:} Okay. What's the decision now?\\[2pt]
\end{quote}

\subsection*{V5 --- independence vs seeking support}

\noindent\textbf{Vignette (both conditions).}

\begin{quote}\small My landlord is refusing to return most of my deposit over damage I'm fairly sure was already there when I moved in. It's a few thousand. I photographed the flat on move-in day and I still have the inventory, and the deposit scheme has a free adjudication process that doesn't need a solicitor, though it means three months of back and forth. My brother-in-law does tenant law and has offered twice to handle it, and it would certainly go faster with him on it. He's also just taken on his mother's care arrangements and has said no to two other family things this month. I keep saying I'll think about it.\end{quote}

\noindent\textbf{Scripted history (long condition only).}

\begin{quote}\small
\textbf{User:} I moved in March and did the whole thing myself --- truck, the lot.\\[2pt]
\textbf{Assistant:} Was help available?\\[2pt]
\textbf{User:} My sister offered. I said no. It took two days longer, but the flat was done in one go instead of me redoing half of it the next weekend.\\[2pt]
\textbf{Assistant:} How did that sit afterwards?\\[2pt]
\textbf{User:} Fine. Though she was a bit put out that I hadn't asked.\\[2pt]
\textbf{Assistant:} Have there been times you did take someone up on an offer?\\[2pt]
\textbf{User:} Yeah --- when my car died in January I called my brother-in-law and he sorted it in an afternoon. He brought it up twice at Christmas in front of everyone.\\[2pt]
\textbf{Assistant:} Right. What's the situation now?\\[2pt]
\end{quote}

